\documentclass[11pt]{article}
\usepackage[margin=1in]{geometry}
\usepackage{times}
\usepackage{hyperref}
\hypersetup{colorlinks=true, linkcolor=black, urlcolor=blue, citecolor=black}
\title{Why We Can't Walk Through Walls}
\author{Flyxion \\ \small Independent Researcher}
\date{}
\begin{document}
\maketitle

\begin{abstract}
Claims that human beings might learn to pass through solid matter, whether framed as a meditative "quantum jumping" skill, a ghost story convention, or a loose misreading of quantum tunneling, share a common evidentiary and mechanistic gap with the antigravity claims surveyed elsewhere in this series: an appeal to a real physical phenomenon, quantum tunneling, borrowed from a regime many orders of magnitude removed from the one in which the claim is actually made. This paper reviews why ordinary solid matter resists interpenetration, an outcome of electromagnetic repulsion between electron clouds reinforced by the Pauli exclusion principle rather than of any adjustable or trainable property of the observer, and computes why the tunneling probability for a macroscopic body through a macroscopic barrier is not merely small but is smaller than any number that has ever mattered to any process in the observable universe.
\end{abstract}

\section{Introduction}

The claim that a person might, through training, belief, or a specific meditative technique, learn to walk through a solid wall recurs across a range of otherwise unrelated contexts, from folk ghost lore to some contemporary self-improvement programs invoking quantum mechanics, and from science fiction convention to informal misreadings of quantum tunneling shared in casual conversation. Unlike most of the claims surveyed elsewhere in this series, this one usually does not even propose a specific mechanism; it more often gestures at "quantum" phenomena in general as license for the possibility, without specifying which quantum effect is meant to be doing the work, or at what scale.

\section{Why Solid Matter Resists Interpenetration in the First Place}

Ordinary solid objects, a hand and a wall alike, are composed overwhelmingly of empty space at the atomic scale, a fact frequently and correctly cited by popularizers of this claim as its starting point. What that framing omits is that the resistance two solid objects offer to interpenetration has never depended on how much of that space is filled with nucleus versus vacuum; it depends on the electromagnetic force between the electron clouds of the atoms at each surface, which repel one another strongly at short range, and on the Pauli exclusion principle, which forbids the wall's electrons and the hand's electrons from occupying the same quantum states, a constraint with no classical analogue and no dependence on the observer's mental state, training, or belief. This is not a matter of atoms being too crowded to pass through in a loose, spatial sense; it is a matter of two independently well confirmed physical principles, electromagnetic repulsion and quantum exclusion, jointly producing an effective, extremely stiff contact force at the boundary between any two ordinary solid objects, a force with the same physical status as the one preventing a table from sinking through the floor.

\section{What Quantum Tunneling Actually Permits, and at What Scale}

Quantum tunneling, the genuine phenomenon most often invoked to lend plausibility to walking-through-walls claims, describes the nonzero probability that a quantum particle will be found on the far side of a potential energy barrier it classically lacks the energy to cross, a real and experimentally confirmed effect exploited in scanning tunneling microscopy, nuclear alpha decay, and semiconductor device physics. The tunneling probability depends exponentially on the barrier's width, height, and the mass of the tunneling object, meaning it falls off catastrophically as any of these three quantities increases, and a human body attempting to tunnel through a wall involves a barrier many meters in effective width by atomic standards and an object whose mass exceeds that of an individual tunneling particle, an electron or alpha particle, by on the order of thirty orders of magnitude, an exponent that appears inside the exponential itself, driving the resulting probability not merely toward zero but toward a number with more zeroes after the decimal point than there are particles in the observable universe, a quantity for which "vanishingly small" is itself a considerable overstatement of how frequently the event would be expected to occur across the entire remaining lifetime of the universe.

\section{Why "Quantum" Does Not License the Extrapolation}

The rhetorical move common to this claim, and shared with several devices surveyed elsewhere in this series, is to correctly cite a real quantum phenomenon operating in one regime, single subatomic particles and barriers on the scale of nanometers, and then informally extend it, with no accompanying calculation, to an entirely different regime, macroscopic bodies and barriers on the scale of meters, where the same underlying mathematics that makes tunneling occasionally observable at the small scale makes it, at the human scale, indistinguishable from impossible for any practical purpose, a distinction the phrase "quantum mechanics says anything is possible" consistently elides.

\section{The Entropic Gravity Framing}

This essay's subject concerns electromagnetic repulsion and quantum tunneling rather than gravitation, and is included in this series less as a gravitational case and more as a companion piece illustrating the same evidentiary pattern this series has traced throughout: a genuine physical phenomenon, real and well confirmed in its native regime, extended by informal analogy into a regime many orders of magnitude removed from where its effects are actually measurable, without the calculation that would reveal the extension's scale is the entire content of the objection.

\section{Objections}

Proponents of meditative or belief-based approaches to this claim generally do not offer a specific physical mechanism to critique, relying instead on anecdote and subjective experience of altered perception during meditative states, which is a genuine and separately studied phenomenon in contemplative neuroscience, but altered subjective perception of one's body or surroundings during a meditative state has no established causal relationship to the electromagnetic and quantum-mechanical forces actually responsible for solid matter's resistance to interpenetration, a gap this specific version of the claim does not attempt to close with a mechanism at all.

\section{Conclusion}

Solid matter resists interpenetration because of electromagnetic repulsion and the Pauli exclusion principle acting between electron clouds, forces with no dependence on belief, training, or observer state, and the quantum tunneling effect genuinely available to subatomic particles at nanometer barriers becomes, for a macroscopic human body at a wall several meters wide, a probability whose exponent alone renders the event outside the range of anything that has ever happened or will happen in the observable universe.

\end{document}
